problem:  
Let \( (M^n, g) \) be a smooth closed manifold with the rational homology of \( S^n \). For each conformal class \([g]\), consider the set  
\[
\mathcal{S}(M,[g]) = \{ f \in C^\infty(M) \mid \text{there exists } \tilde{g} \in [g] \text{ such that } R_{\tilde{g}} = f \}.
\]  
It is known that for \( M = S^n \), the solvability of the prescribed scalar curvature equation imposes sharp conditions on \( f \) (e.g. the Kazdan–Warner obstruction).   
  
**Question:** Suppose \(M^n\) is an *exotic sphere* (\(n \ge 7\)). Does there exist a conformal class \([g]\) on \(M\) such that \(\mathcal{S}(M,[g]) = C^\infty(M)\), i.e. every smooth function is realizable as the scalar curvature of some metric conformal to \(g\)?  
In other words, can a non-standard differential structure on a sphere remove all global obstructions to prescribing scalar curvature in its conformal classes?  

Why is it a "good" problem:  
This problem explores whether the smooth structure of a sphere influences the analytic solvability of the prescribed scalar curvature equation. It straddles the interface of **differential topology** (via exotic smooth structures) and **geometric analysis** (via the Yamabe-type PDE and scalar curvature prescription). Resolving it would illuminate whether scalar curvature, an analytic object, is sensitive to exotic smoothness—a question at the heart of modern differential geometry. The formulation is simple yet profound: it asks if the analytic flexibility of scalar curvature can detect the difference between the standard and exotic spheres. A positive or negative resolution could lead to a new geometric invariant distinguishing smooth structures, or conversely suggest analytic equivalence across them, either outcome representing a major advance.